\PassOptionsToPackage{unicode=true}{hyperref} % options for packages loaded elsewhere
\PassOptionsToPackage{hyphens}{url}
%
\documentclass[14pt,]{extarticle}
\usepackage{lmodern}
\usepackage{amssymb,amsmath}
\usepackage{ifxetex,ifluatex}
\usepackage{fixltx2e} % provides \textsubscript
\ifnum 0\ifxetex 1\fi\ifluatex 1\fi=0 % if pdftex
  \usepackage[T1]{fontenc}
  \usepackage[utf8]{inputenc}
  \usepackage{textcomp} % provides euro and other symbols
\else % if luatex or xelatex
  \usepackage{unicode-math}
  \defaultfontfeatures{Ligatures=TeX,Scale=MatchLowercase}
    \setmainfont[]{Times New Roman}
  \ifxetex
    \usepackage{xeCJK}
    \setCJKmainfont[]{Microsoft JhengHei}
  \fi
  \ifluatex
    \usepackage[]{luatexja-fontspec}
    \setmainjfont[]{Microsoft JhengHei}
  \fi
\fi
% use upquote if available, for straight quotes in verbatim environments
\IfFileExists{upquote.sty}{\usepackage{upquote}}{}
% use microtype if available
\IfFileExists{microtype.sty}{%
\usepackage[]{microtype}
\UseMicrotypeSet[protrusion]{basicmath} % disable protrusion for tt fonts
}{}
\IfFileExists{parskip.sty}{%
\usepackage{parskip}
}{% else
\setlength{\parindent}{0pt}
\setlength{\parskip}{6pt plus 2pt minus 1pt}
}
\usepackage{hyperref}
\hypersetup{
            pdfborder={0 0 0},
            breaklinks=true}
\urlstyle{same}  % don't use monospace font for urls
\usepackage[margin=1in]{geometry}
\usepackage{longtable,booktabs}
% Fix footnotes in tables (requires footnote package)
\IfFileExists{footnote.sty}{\usepackage{footnote}\makesavenoteenv{longtable}}{}
\setlength{\emergencystretch}{3em}  % prevent overfull lines
\providecommand{\tightlist}{%
  \setlength{\itemsep}{0pt}\setlength{\parskip}{0pt}}
\setcounter{secnumdepth}{0}
% Redefines (sub)paragraphs to behave more like sections
\ifx\paragraph\undefined\else
\let\oldparagraph\paragraph
\renewcommand{\paragraph}[1]{\oldparagraph{#1}\mbox{}}
\fi
\ifx\subparagraph\undefined\else
\let\oldsubparagraph\subparagraph
\renewcommand{\subparagraph}[1]{\oldsubparagraph{#1}\mbox{}}
\fi

% set default figure placement to htbp
\makeatletter
\def\fps@figure{htbp}
\makeatother


\date{}

\begin{document}

\hypertarget{week-05-ux4f5cux696dux56deux994b}{%
\section{Week 05 作業回饋}\label{week-05-ux4f5cux696dux56deux994b}}

\begin{quote}
\textbf{評分標準：} 依據 \texttt{week05\_rubric\_v2.xlsx} 之 13
項評分指標，滿分 100 分。 各大項配分：A. 實驗架構（Coder）50 分 ／ B.
Builder 版本 20 分 ／ C. 程式品質 15 分 ／ D. 文件與繳交 15 分
\end{quote}

\begin{center}\rule{0.5\linewidth}{0.4pt}\end{center}

\hypertarget{ux5433ux5fc3ux5713-96100}{%
\subsection{113892001 吳心圓 ---
96／100}\label{ux5433ux5fc3ux5713-96100}}

\hypertarget{ux7e3dux8a55}{%
\subsubsection{總評}\label{ux7e3dux8a55}}

刺激設計具獨到巧思------target 為垂直翻轉（y 取反）、distractor
為水平翻轉（x
取反），使二者視覺差異更具挑戰性，符合知覺辨識實驗的科學考量。距離門檻過濾（\texttt{np.linalg.norm\ \textgreater{}\ 5px}）是很好的效能優化選擇。小幅扣分在於
Escape 偵測不完整及未附範例 CSV 檔案。

\hypertarget{ux5404ux9805ux8a55ux8a9e}{%
\subsubsection{各項評語}\label{ux5404ux9805ux8a55ux8a9e}}

\begin{longtable}[]{@{}lll@{}}
\toprule
\begin{minipage}[b]{0.09\columnwidth}\raggedright
指標\strut
\end{minipage} & \begin{minipage}[b]{0.03\columnwidth}\raggedright
得分\strut
\end{minipage} & \begin{minipage}[b]{0.80\columnwidth}\raggedright
評語\strut
\end{minipage}\tabularnewline
\midrule
\endhead
\begin{minipage}[t]{0.09\columnwidth}\raggedright
A1 試驗迴圈\strut
\end{minipage} & \begin{minipage}[t]{0.03\columnwidth}\raggedright
15/15\strut
\end{minipage} & \begin{minipage}[t]{0.80\columnwidth}\raggedright
\texttt{for\ trial\ in\ range(n\_trials)} 正確執行 10 輪，以
\texttt{n\_trials} 變數控制，易於調整。\strut
\end{minipage}\tabularnewline
\begin{minipage}[t]{0.09\columnwidth}\raggedright
A2 刺激呈現與計時\strut
\end{minipage} & \begin{minipage}[t]{0.03\columnwidth}\raggedright
10/10\strut
\end{minipage} & \begin{minipage}[t]{0.80\columnwidth}\raggedright
繪圖階段以 \texttt{core.Clock} 控制 10 秒計時；選擇階段分為前 5
秒純觀看（target +
distractor）與後續按鍵收集，將「知覺編碼」與「反應執行」分離，這在認知實驗設計中是有意義的做法，可減少速度---準確度權衡的干擾。\strut
\end{minipage}\tabularnewline
\begin{minipage}[t]{0.09\columnwidth}\raggedright
A3 反應收集與正確性\strut
\end{minipage} & \begin{minipage}[t]{0.03\columnwidth}\raggedright
15/15\strut
\end{minipage} & \begin{minipage}[t]{0.80\columnwidth}\raggedright
以 \texttt{left}/\texttt{right} 方向鍵配合 \texttt{side}
變數判斷正確性，邏輯正確。\strut
\end{minipage}\tabularnewline
\begin{minipage}[t]{0.09\columnwidth}\raggedright
A4 區塊結構\strut
\end{minipage} & \begin{minipage}[t]{0.03\columnwidth}\raggedright
10/10\strut
\end{minipage} & \begin{minipage}[t]{0.80\columnwidth}\raggedright
觀看階段與反應階段的時序分離提供了試驗內的結構化節奏控制。\strut
\end{minipage}\tabularnewline
\begin{minipage}[t]{0.09\columnwidth}\raggedright
B1 Builder 流程重建\strut
\end{minipage} & \begin{minipage}[t]{0.03\columnwidth}\raggedright
10/10\strut
\end{minipage} & \begin{minipage}[t]{0.80\columnwidth}\raggedright
Builder
版本完整重現實驗流程，且同樣實作了先觀看再按鍵的兩階段設計。\strut
\end{minipage}\tabularnewline
\begin{minipage}[t]{0.09\columnwidth}\raggedright
B2 Builder 資料紀錄\strut
\end{minipage} & \begin{minipage}[t]{0.03\columnwidth}\raggedright
5/5\strut
\end{minipage} & \begin{minipage}[t]{0.80\columnwidth}\raggedright
使用
\texttt{thisExp.addData(\textquotesingle{}is\_correct\textquotesingle{},\ int(is\_correct))}
與
\texttt{thisExp.addData(\textquotesingle{}side\_of\_target\textquotesingle{},\ side)}
將結果正確記入資料檔。\strut
\end{minipage}\tabularnewline
\begin{minipage}[t]{0.09\columnwidth}\raggedright
B3 Code Component / JS\strut
\end{minipage} & \begin{minipage}[t]{0.03\columnwidth}\raggedright
5/5\strut
\end{minipage} & \begin{minipage}[t]{0.80\columnwidth}\raggedright
Code Component 邏輯正確，與 Coder 版本一致。\strut
\end{minipage}\tabularnewline
\begin{minipage}[t]{0.09\columnwidth}\raggedright
C1 可讀性與組織\strut
\end{minipage} & \begin{minipage}[t]{0.03\columnwidth}\raggedright
5/5\strut
\end{minipage} & \begin{minipage}[t]{0.80\columnwidth}\raggedright
程式碼簡潔，使用 \texttt{np.linalg.norm}
做距離門檻過濾是有效率的寫法，變數命名清楚（\texttt{drawing\_points}、\texttt{flipped\_points}、\texttt{distractor\_points}）。\strut
\end{minipage}\tabularnewline
\begin{minipage}[t]{0.09\columnwidth}\raggedright
C2 Escape 處理\strut
\end{minipage} & \begin{minipage}[t]{0.03\columnwidth}\raggedright
\textbf{3/5}\strut
\end{minipage} & \begin{minipage}[t]{0.80\columnwidth}\raggedright
⚠️ 繪圖階段有偵測
Escape（\texttt{if\ \textquotesingle{}escape\textquotesingle{}\ in\ event.getKeys():\ win.close();\ core.quit()}），但選擇階段的
\texttt{while\ not\ responded} 迴圈中未包含 Escape
偵測。若受試者在選擇階段想退出，需等待該輪結束或強制關閉程式。\textbf{（-2）}\strut
\end{minipage}\tabularnewline
\begin{minipage}[t]{0.09\columnwidth}\raggedright
C3 邊界情況處理\strut
\end{minipage} & \begin{minipage}[t]{0.03\columnwidth}\raggedright
5/5\strut
\end{minipage} & \begin{minipage}[t]{0.80\columnwidth}\raggedright
當繪圖為空時以 \texttt{{[}(0,0),\ (10,10){]}} 預設值填充，避免 ShapeStim
崩潰。\strut
\end{minipage}\tabularnewline
\begin{minipage}[t]{0.09\columnwidth}\raggedright
D1 繳交完整度\strut
\end{minipage} & \begin{minipage}[t]{0.03\columnwidth}\raggedright
5/5\strut
\end{minipage} & \begin{minipage}[t]{0.80\columnwidth}\raggedright
同時繳交 \texttt{.py} 與 \texttt{.psyexp}，附
\texttt{content.html}。\strut
\end{minipage}\tabularnewline
\begin{minipage}[t]{0.09\columnwidth}\raggedright
D2 範例輸出\strut
\end{minipage} & \begin{minipage}[t]{0.03\columnwidth}\raggedright
\textbf{3/5}\strut
\end{minipage} & \begin{minipage}[t]{0.80\columnwidth}\raggedright
⚠️ 未附上範例 CSV 檔案。Builder 版本透過 \texttt{thisExp.addData}
可自動儲存，但繳交時未包含實際執行後的範例資料檔。\textbf{（-2）}\strut
\end{minipage}\tabularnewline
\begin{minipage}[t]{0.09\columnwidth}\raggedright
D3 科學設計品質\strut
\end{minipage} & \begin{minipage}[t]{0.03\columnwidth}\raggedright
5/5\strut
\end{minipage} & \begin{minipage}[t]{0.80\columnwidth}\raggedright
垂直翻轉作為 target、水平翻轉作為 distractor
的設計增加辨識難度，符合知覺辨識研究中對干擾項設計的考量。\strut
\end{minipage}\tabularnewline
\bottomrule
\end{longtable}

\hypertarget{ux6539ux9032ux5efaux8b70}{%
\subsubsection{改進建議}\label{ux6539ux9032ux5efaux8b70}}

\begin{enumerate}
\def\labelenumi{\arabic{enumi}.}
\tightlist
\item
  \textbf{補充 Escape 偵測（+2 分）：} 在 \texttt{while\ not\ responded}
  迴圈中加入
  \texttt{if\ \textquotesingle{}escape\textquotesingle{}\ in\ event.getKeys():\ win.close();\ core.quit()}
  檢查，確保受試者可隨時退出。
\item
  \textbf{附上範例輸出（+2 分）：} 執行 Builder 版本一次，將產生的 CSV
  檔案一併繳交，以利驗證實驗功能。
\end{enumerate}

\end{document}
